Este proyecto busca desarrollar una plataforma web capaz de administrar, publicar, participar
y darle seguimiento a proyectos de extensión de la Universidad del Valle de Guatemala.
Para lograr esto, la plataforma web constará de dos partes principales: un \textit{Backend} 
y un \textit{Frontend}.

\section{Entregables}
\subsection{\textit{Backend}}
Como entregable del \textit{Backend} se tendrá una API con un conjunto de \textit{endpoints}
que permitan realizar operaciones según el rol del usuario.

\subsubsection{Todos los usuarios}
\begin{itemize}
    \item Registro de usuario
    \item Inicio de sesión
    \item Verificación de correo electrónico
    \item Restablecimiento de contraseña
    \item Actualización de información personal y cambio de contraseña
    \item Búsqueda y visualización de proyectos
    \item Visualización de motivaciones
\end{itemize}

\subsubsection{Estudiantes}
\begin{itemize}
    \item Postulación a proyectos
    \item Participación en un proyecto a la vez
    \item Abandonar un proyecto
    \item Registro de horas trabajadas en un proyecto, incluyendo evidencias
    \item Visualización de horas aprobadas en proyectos pasados
    \item Visualización de historial de postulaciones y horas
\end{itemize}

\subsubsection{Organizaciones externas y directores}
\begin{itemize}
    \item Creación, edición y eliminación de proyectos con archivos (recursos) ligados
    \item Iniciar, finalizar o cancelar un proyecto propio
    \item Aceptar o rechazar postulaciones a proyectos propios
    \item Visualización de participantes en proyectos propios
    \item Eliminación de participantes en proyectos propios
    \item Aprobación o rechazo de horas registradas por estudiantes en proyectos propios
\end{itemize}

\subsubsection{Directores}
\begin{itemize}
    \item Aprobación o rechazo de proyectos
    \item Aprobación o rechazo de cuentas de organizaciones externas
\end{itemize}

\subsection{\textit{Frontend}}
Como entregable del \textit{Frontend} se tendrá una interfaz gráfica alineada a los diseños
previamente construidos y validados. Esta interfaz deberá de permitir a los usuarios finales
realizar las operaciones descritas anteriormente según su rol.

Las pantallas para la plataforma son las siguientes:

\begin{itemize}
    \item Home
    \item Registro
    \item Inicio de sesión
    \item Verificación de correo electrónico
    \item Restablecimiento de contraseña
    \item Perfil
    \item Edición de perfil
    \item Cambio de contraseña
    \item Proyectos propios
    \item Creación de proyecto
    \item Edición de proyecto
    \item Búsqueda de proyectos
    \item Visualización de proyecto
    \item Configuración de proyecto
    \item Proyecto actual
    \item Historial de proyectos
    \item Registro de horas
    \item Motivaciones
    \item Proyectos pendientes de aprobación
    \item Organizaciones externas pendientes de aprobación
\end{itemize}

\section{Límites}
Los siguientes puntos no están incluidos dentro del alcance de este proyecto:
\begin{itemize}
    \item Creación de diseños para las pantallas y experiencia de usuario.
    \item Pruebas de tiempo de carga y rendimiento para el \textit{Frontend}.
    \item Implementación de funcionalidades adicionales no especificadas en los entregables.
\end{itemize}
